\label{ch:discussion}

\section{Overview}
\label{sec:disc-overview}

Chapter~\ref{ch:results} answered RQ1--RQ3 for each of three generators. This chapter interprets those answers against the research gap (Chapter~\ref{ch:research-gap}) and the literature (Chapter~\ref{ch:lit-review}), states what the cross-generator design does and does not establish, and sets out the thesis's limitations honestly before Chapter~\ref{ch:conclusion} states its contributions.

\section{What the Results Say About Each Research Question}
\label{sec:disc-rqs}

\subsection{RQ1: hallucination is real, but its shape depends on the generator}
\label{sec:disc-rq1}

Baseline paired HR is high in every generator tested (57.5--82.8\%), consistent with the magnitude Dahl et al.\ and Magesh et al.\ report for legal QA and commercial legal-research tools (58--88\% and 17--33\% respectively) \cite{dahl2024largelegalfictions,magesh2025hallucinationfree}. That consistency matters: it is the first evidence in this thesis that the citation-hallucination problem those studies document in QA and closed-tool settings reproduces in open, long-document \emph{legal summarization}, the setting Chapter~\ref{ch:research-gap}'s second consolidated gap identifies as absent from the literature. Dahl et al.\ and Magesh et al.\ do not summarize 75,000-word civil-rights dockets; this thesis's baseline numbers are the first evidence that the same failure mode appears there too, at comparable severity.

But RQ1's more specific finding, that citation-attempt density and failure \emph{type} both scale with generator choice, not with the underlying task (Section~\ref{sec:results-rq1}), was not something the literature review's gap analysis anticipated, because no reviewed work varies the generator architecture while holding the verification pipeline fixed. This is a genuine finding of this thesis's cross-model design, not a confirmation of prior work: PRIMERA's near-absence of citation attempts (0.29/case) could plausibly have been read as a property of legal summarization under heavy truncation, and it took building two further generators specifically to rule that out (Section~\ref{sec:module-a-crossmodel}).

\subsection{RQ2: verification-plus-correction closes the gap detection-only methods leave open}
\label{sec:disc-rq2}

Chapter~\ref{ch:research-gap}'s first and third consolidated gaps state that retrieval reduces but does not eliminate hallucination, and that detection methods (HalluGraph, SelfCheckGPT, TRUE) flag risk without repairing the output \cite{lewis2020rag,noel2025hallugraph,manakul2023selfcheckgpt,honovich2022true}. RQ2's results speak directly to both. On the first: this thesis's own retrieval variant (LLM~+~retrieval) still starts at 70.5\% baseline HR, exactly the double-digit residual hallucination rate retrieval alone would leave, matching what Magesh et al.\ report for commercial RAG-based tools. Symbolic verification and correction is what takes it to 24.2\%, not retrieval by itself. The third gap holds up too. This thesis's pipeline does not stop at a HalluGraph- or SelfCheckGPT-style risk score. Every citation Module~C can resolve with a safe identity match is rewritten in place (fix rates 31.0--64.9\% across generators), not merely flagged, and the correction is auditable per citation (Section~\ref{sec:module-c}) rather than a black-box confidence number.

The result the literature gives no basis to predict is the zero-regression finding: across 1{,}083 paired citations and three architecturally different generators, correction never introduced a new hallucination (Table~\ref{tab:rq2-discordant}). Detection-only baselines cannot produce this kind of evidence at all, because they never edit the output; the comparison is not detection accuracy versus correction accuracy but whether an editing system can be trusted not to make things worse, which is the practical question a professional legal workflow actually asks. The consistency of that zero across a 24$\times$ range in citation volume is what makes it a property of the correction \emph{policy}: specifically, its refusal to substitute a contradicting case name (Section~\ref{sec:disc-rq2-why}) and its use of NotFound/Malformed as unconditional redact signals, rather than an artefact of one favourable generator.

\subsection{RQ3: the quality-faithfulness trade-off this pipeline was designed to avoid does not appear}
\label{sec:disc-rq3}

Chapter~\ref{ch:research-gap}'s first consolidated gap is that benchmark quality and professional reliability are not the same thing: Multi-LexSum's own PRIMERA baseline pairs competitive ROUGE with a 0.43/3.0 human faithfulness rating \cite{shen2022multilexsum}, and Maynez et al.\ and Pagnoni et al.\ document the same decoupling in general-domain summarization \cite{maynez2020faithfulness,pagnoni2021factuality}. A pipeline that repaired citations by aggressively rewriting surrounding text could close the faithfulness gap while reopening a quality gap of its own. RQ3's results (Table~\ref{tab:rq3-results}) show this pipeline does not do that: ROUGE-L\textsubscript{sum} moves by at most 0.07 percentage points in any generator, against a $\pm2$pp tolerance, and BERTScore F1 moves by at most 0.0003. The observed deltas are one to two orders of magnitude smaller than the tolerance itself, not a coincidental win against a lenient target.

The reason is architectural rather than incidental: symbolic correction (Section~\ref{sec:correction}) edits only the character spans eyecite identifies as citations, typically a few dozen characters inside a summary of several hundred words, and never rewrites surrounding prose. The optional LLM-rewrite condition, which does touch surrounding text, is deliberately excluded from the RQ3 headline comparison (Section~\ref{sec:correction}, ``Optional LLM rewrite'') precisely because it is the ablation most likely to reproduce the quality cost RQ3 asks whether this pipeline avoids.

\subsection{Why the correction policy's conservatism is doing the work}
\label{sec:disc-rq2-why}

Chapter~\ref{ch:research-gap}'s fourth gap, unsafe repair without identity checks, where a mismatched citation is silently replaced by an unrelated real case, is the specific failure mode Section~\ref{sec:method-metrics}'s Ambiguous label and re-anchor policy (Section~\ref{sec:module-c}, ``Why Ambiguous is not `Mismatch repair'\,'') were designed against. The zero-regression result in Section~\ref{sec:disc-rq2} is the empirical confirmation that this design choice matters. A pipeline that resolved every Mismatch by substituting CourtListener's returned case name would repair more citations per case, but nothing in this evaluation would have detected whether some of those repairs inserted the wrong case, since draft-vs-final HR alone cannot distinguish ``still hallucinated'' from ``hallucinated differently.'' The discordant-pair breakdown in Table~\ref{tab:rq2-discordant} is the check that answers this directly, and it answers it cleanly across all three generators.

\section{Generalization Across Generators: What the Cross-Model Design Establishes}
\label{sec:disc-generalization}

The cross-generator comparison (Section~\ref{sec:results-crossgen}) was motivated by a single observation during development, PRIMERA's implausibly low citation-attempt count, but its result is broader than resolving that one question. Module~C and the correction policy were built and tuned against PRIMERA's citation profile (few attempts, heavily NotFound) before the LLM-based generators existed. Final HR nonetheless lands in a narrow 22.2--27.6\% band for all three, despite a 25-point spread in baseline HR and a 24$\times$ spread in citation volume. A pipeline whose correction logic happened to exploit some quirk of PRIMERA's output would be expected to perform unevenly once that quirk disappeared; instead, performance is remarkably stable. This is the strongest evidence available in this thesis that the taxonomy and correction policy (Chapter~\ref{ch:methodology}) are addressing a general property of citation verification against an external database, not overfitting to one checkpoint's failure modes.

The qualitative finding in Section~\ref{sec:results-qualitative}, that citation-density retrieval suppresses degenerate same-case citation fabrication from 13.7\% to 0.7\% of citations, adds a second, independent line of evidence for retrieval's value beyond the aggregate HR numbers, and one the literature's retrieval-versus-verification framing (Chapter~\ref{ch:research-gap}, gap 1; Chapter~\ref{ch:lit-review}, Theme 2) does not distinguish: retrieval here changes \emph{which kind} of wrong citations are, from unconstrained fabrication toward identifiable near-misses that the taxonomy can label and the correction policy can act on, and not only \emph{how often} they are wrong.

\section{Limitations}
\label{sec:disc-limitations}

\paragraph{Human validation is not yet part of this evaluation.} Every number in Chapter~\ref{ch:results} comes from the automated Module~C taxonomy and automated quality metrics. Two pieces of human-in-the-loop evaluation remain open: (i)~a manual audit of the automated Verified/Mismatch/Ambiguous/NotFound/Malformed labels against human legal judgment on the 200-case stratified sample already exported (Section~\ref{sec:audit-subset}), which would establish how much the automated taxonomy itself over- or under-counts hallucination relative to a human reviewer; and (ii)~a human readability evaluation of corrected summaries, since redaction in particular can leave a grammatically abrupt gap where a citation was removed, and no reader-facing measurement of that cost exists yet. Both are treated as immediate next steps (Chapter~\ref{ch:conclusion}) rather than completed validation, and every claim in this thesis about hallucination rates and quality preservation should be read as automated-metric evidence pending that human check.

\paragraph{The evaluation set is a deliberate stress test, not a random sample.} The 100 cases used throughout Chapter~\ref{ch:results} are the most citation-dense cases in the test split (Section~\ref{sec:eval-case-selection}), chosen because paired HR is uninformative on cases with few or no citations. This is the right design for answering RQ1--RQ3 with statistical power, but it means the reported HR values describe pipeline behaviour where citations are common and hallucination is most consequential, not an unconditional average over the full test split; a case with zero gold citations contributes nothing to this evaluation and may behave differently.

\paragraph{CourtListener is treated as ground truth, with the limitations that implies.} Module~C's Verified/Mismatch/Ambiguous/NotFound labels are only as correct as CourtListener's own database coverage and metadata accuracy (Section~\ref{sec:data-courtlistener}). A citation to a real case CourtListener has not indexed would be labelled NotFound and redacted exactly as a fabricated citation would be, understating precision in a way this evaluation cannot currently distinguish without the manual audit above.

\paragraph{The degenerate-repetition failure mode is suppressed, not eliminated.} Section~\ref{sec:results-qualitative} shows citation-density retrieval reduces fabricated same-case citation runs from 13.7\% to 0.7\% of citations, and the correction policy redacts or flags every individual instance correctly. But 0.7\% is not zero, and the mitigation is a side effect of a design decision (source ranking) made for a different reason (citation coverage), not a purpose-built defence against this specific pattern.

\paragraph{Truncation remains a hard constraint on what each generator can see.} PRIMERA's 4{,}096-token ceiling and the LLM-based generators' 8{,}192-token budget (Section~\ref{sec:module-a-crossmodel}) are both far below the corpus's median source length of 24{,}546 words (Chapter~\ref{ch:eda}). Every generator in this evaluation drafts from a truncated view of the docket; symbolic verification checks whether cited authorities are real, not whether the summary omits material the generator never saw.

\paragraph{Scope is Document$\rightarrow$Long, U.S.\ civil-rights litigation, English.} Short/tiny granularities, non-English proceedings, and case-law systems outside CourtListener's coverage are outside this evaluation (Section~\ref{sec:eval-case-selection}); the taxonomy and correction policy's generalization claim (Section~\ref{sec:disc-generalization}) is a claim about generator architecture, evaluated within this single corpus and jurisdiction.

\section{Threats to Validity}
\label{sec:disc-threats}

\textbf{Construct validity.} Paired HR (Section~\ref{sec:method-metrics}) counts an Ambiguous citation left unresolved as \texttt{still\_bad}, the same outcome as a confirmed fabrication. This is a deliberate conservative choice: the alternative would let the metric reward leaving a genuinely uncertain citation untouched. But it means paired HR is better read as ``fraction of citations this pipeline could not certify as safe'' than as a pure fabrication rate.

\textbf{Internal validity.} All three generators share Modules~B and~C and the correction policy; a latent bug in extraction or verification would affect every generator identically and could not be detected by the cross-generator agreement argued for in Section~\ref{sec:disc-generalization}. The regression-tested implementation (Chapter~\ref{ch:methodology}) reduces but does not eliminate this risk.

\textbf{Statistical power.} McNemar tests in Chapter~\ref{ch:results} are computed over citation counts (29 to 689), not case counts (100 in every generator); PRIMERA's $n=29$ paired citations is the thinnest of the three, though still sufficient for $p=3.05\times10^{-5}$.

\section{Chapter Summary}
\label{sec:disc-summary}

The results in Chapter~\ref{ch:results} confirm the literature's core premises: legal citation hallucination is frequent and severe, retrieval alone does not resolve it, and quality metrics can be preserved without sacrificing faithfulness gains, while adding two findings the reviewed literature does not establish: that a symbolic correction policy tuned against one generator's failure profile generalizes to generators with a 24$\times$ different citation volume without HR degrading, and that it does so with a measured zero rate of introduced regressions across every generator tested. The main open validation step is human-in-the-loop confirmation that the automated taxonomy and the readability of corrected output agree with human judgment, explicitly scoped as future work in Chapter~\ref{ch:conclusion}, not as a caveat that undermines the automated results reported here.
